% ============================================================
%  Section 4 — Data
% ============================================================
\section{Synthetic Corpus}
\label{sec:data}

\subsection{Document Generation}

We generate 300 documents (150 per company) using a deterministic Python generator
seeded at 42, so every run produces the identical corpus. Documents are PDFs with
a text layer (ReportLab~\cite{reportlab}) plus a degraded photo variant.
All text is in French.

Each generated document contains a selected set of accounting-style fields:
supplier name, a synthetic matricule fiscal (7 digits + letter), date in
\textit{jj/mm/aaaa} format,
line items with HTVA/TVA/TTC amounts, a Timbre Fiscal line, and the document total.
Amounts use the French format with a narrow no-break space as thousands separator
(e.g., \textit{1\,011,500\,TND}).

\textbf{Company A} uses TND currency, 19\% VAT on all lines, and six-character
account codes such as \textit{606100}.

\textbf{Company B} uses EUR and TND, tiered VAT at 7\%, 13\%, or 19\% per line,
and account codes of the form \textit{6060-EN001}.

Photo variants add synthetic degradation: rotation between 8° and 15°,
Gaussian noise with $\sigma \in [15,20]$, and JPEG quality between 55 and 65.
They are not a substitute for independently collected, real-world scans.

\subsection{Error Injection}

30\% of documents contain exactly one injected error:

\begin{itemize}[leftmargin=*, itemsep=2pt]
  \item \textbf{Arithmetic mismatch}: the TTC printed on the document does not
        equal HTVA + TVA, triggering the Validator's deterministic check.
  \item \textbf{Missing matricule}: the supplier's tax identifier field is blank.
  \item \textbf{Duplicate ID}: the document is marked as a copy of an existing
        one via a \texttt{duplicate\_of} field.
\end{itemize}

Documents with injected errors should exit as \texttt{ESCALATED}; clean documents
should exit as \texttt{ACCEPTED}. This gives us a ground-truth label for measuring
Validator F-score.

\subsection{Hidden House Rules}

Each company has four hidden rules stored separately from the public RAG corpus
and represented in generated gold metadata. Company B's rules cover EUR invoice
conversion, review of high-value TechVision invoices, a fixed AssurPlus code,
and a flag for Station Total purchase orders. The corpus generator also assigns
TechVision consulting and audit line items to \texttt{6226-HO001}; seven such
TRAIN examples support the v003 skill exception evaluated in
\cref{sec:skillopt}. These company policies are synthetic and are not official
Tunisian accounting guidance.

\subsection{Splits and Integrity}

\begin{table}[ht]
\caption{Corpus statistics}
\label{tab:corpus}
\centering\scriptsize\renewcommand{\arraystretch}{1.2}
\begin{tabular}{@{}lrrrr@{}}
\toprule
\textbf{Split} & \textbf{Size} & \textbf{Co.\,A} & \textbf{Co.\,B} & \textbf{Purpose} \\
\midrule
TRAIN     & 140 & 70 & 70 & Skill development \\
SELECTION & 60  & 30 & 30 & Baseline eval, validation \\
TEST      & 100 & 50 & 50 & Final comparison (held-out) \\
\midrule
Total & 300 & 150 & 150 & \\
\bottomrule
\end{tabular}
\end{table}

The SkillOpt pilot used a balanced sample of 12 Company B TRAIN and 12 Company B
SELECTION items from the project data, rather than the full split. It did not
evaluate the TEST partition.

Seven automated integrity checks verify the corpus on every generation:
unique IDs across all 300 documents, no document in more than one split,
arithmetic consistency of all gold labels except injected-error documents,
all account codes present in the company chart, house-rule text absent from
public files, all matricule fiscal values matching the 7-digit+letter pattern,
and no generic placeholder text.
